\chapter{Kubity i proste układy kwantowe}

\begin{summarybox}
Kubit jest najprostszym nietrywialnym układem kwantowym. Ma dwuwymiarową
przestrzeń stanów, ale już zawiera najważniejsze idee: superpozycję, amplitudy
zespolone, pomiar probabilistyczny, operatory jako macierze, iloczyn tensorowy i
splątanie. Ten rozdział wprowadza kubity jako praktyczne laboratorium formalizmu
Diraca.
\end{summarybox}

\section{Kubit jako najprostszy układ kwantowy}

Klasyczny bit ma dwie możliwe wartości: \(0\) albo \(1\). Kubit ma dwie stany bazowe, ale jego ogólny stan może być ich superpozycją. Oznaczamy bazę standardową jako
\begin{equation}
    |0\rangle,
    \qquad
    |1\rangle.
\end{equation}
Ogólny stan kubitu ma postać
\begin{equation}
    |\psi\rangle=\alpha|0\rangle+\beta|1\rangle,
\end{equation}
gdzie \(\alpha\) i \(\beta\) są liczbami zespolonymi spełniającymi warunek normalizacji
\begin{equation}
    |\alpha|^2+|\beta|^2=1.
\end{equation}

To nie oznacza, że kubit „jest trochę zerem i trochę jedynką” w klasycznym sensie. Oznacza, że amplitudy \(\alpha\) i \(\beta\) określają prawdopodobieństwa wyników pomiaru oraz fazę względną, która może wpływać na interferencję.

\section{Stany \texorpdfstring{\(|0\rangle\)}{|0>} i \texorpdfstring{\(|1\rangle\)}{|1>}}

W reprezentacji macierzowej wygodnie zapisać
\begin{equation}
    |0\rangle=\begin{pmatrix}1\\0\end{pmatrix},
    \qquad
    |1\rangle=\begin{pmatrix}0\\1\end{pmatrix}.
\end{equation}
Wtedy ogólny stan kubitu jest kolumną
\begin{equation}
    |\psi\rangle=\begin{pmatrix}\alpha\\\beta\end{pmatrix}.
\end{equation}
Odpowiadający bra ma postać
\begin{equation}
    \langle\psi|=\begin{pmatrix}\alpha^* & \beta^*\end{pmatrix}.
\end{equation}
Normalizacja to zwykły iloczyn skalarny:
\begin{equation}
    \langle\psi|\psi\rangle=|\alpha|^2+|\beta|^2=1.
\end{equation}

\section{Superpozycja}

Superpozycja jest liniową kombinacją stanów. Przykładowo stan
\begin{equation}
    |+\rangle=\frac{1}{\sqrt{2}}\left(|0\rangle+|1\rangle\right)
\end{equation}
ma równe prawdopodobieństwa otrzymania \(0\) i \(1\) przy pomiarze w bazie standardowej. Stan
\begin{equation}
    |-\rangle=\frac{1}{\sqrt{2}}\left(|0\rangle-|1\rangle\right)
\end{equation}
ma te same prawdopodobieństwa w bazie \(|0\rangle,|1\rangle\), ale inną fazę względną. Ta faza jest fizycznie istotna, bo może wpływać na wynik kolejnych operacji.

\begin{remarkbox}
Dwa stany mogą dawać te same prawdopodobieństwa w jednej bazie, ale różnić się fazą. Mechanika kwantowa nie jest teorią samych prawdopodobieństw, lecz amplitud.
\end{remarkbox}

\section{Macierze Pauliego}

Podstawowe operatory działające na kubit to macierze Pauliego:
\begin{equation}
    \sigma_x=\begin{pmatrix}0&1\\1&0\end{pmatrix},
    \qquad
    \sigma_y=\begin{pmatrix}0&-i\\i&0\end{pmatrix},
    \qquad
    \sigma_z=\begin{pmatrix}1&0\\0&-1\end{pmatrix}.
\end{equation}
Macierz \(\sigma_x\) zamienia stany bazowe:
\begin{equation}
    \sigma_x|0\rangle=|1\rangle,
    \qquad
    \sigma_x|1\rangle=|0\rangle.
\end{equation}
Macierz \(\sigma_z\) zostawia \(|0\rangle\) bez zmiany, ale zmienia znak \(|1\rangle\):
\begin{equation}
    \sigma_z|0\rangle=|0\rangle,
    \qquad
    \sigma_z|1\rangle=-|1\rangle.
\end{equation}

Macierze Pauliego spełniają relacje komutacyjne
\begin{equation}
    [\sigma_i,\sigma_j]=2i\sum_k\varepsilon_{ijk}\sigma_k.
\end{equation}
Widać tu tę samą strukturę algebraiczną, która pojawiła się przy momencie pędu.

\section{Bramki kwantowe}

Bramka kwantowa to operator unitarny działający na stan kubitu. Unitarność oznacza zachowanie normy:
\begin{equation}
    U^\dagger U=I.
\end{equation}
Najprostsze bramki to identyczność, bramka \(X\), bramka \(Z\) i bramka Hadamarda. Bramka \(X\) jest po prostu \(\sigma_x\). Bramka Hadamarda ma postać
\begin{equation}
    H=\frac{1}{\sqrt{2}}\begin{pmatrix}1&1\\1&-1\end{pmatrix}.
\end{equation}
Działa ona tak:
\begin{equation}
    H|0\rangle=\frac{|0\rangle+|1\rangle}{\sqrt{2}}=|+\rangle,
\end{equation}
\begin{equation}
    H|1\rangle=\frac{|0\rangle-|1\rangle}{\sqrt{2}}=|-\rangle.
\end{equation}

W Mathematice można to zapisać następująco:
\begin{verbatim}
ket0 = {1, 0};
ket1 = {0, 1};
H = 1/Sqrt[2] {{1, 1}, {1, -1}};
H.ket0
H.ket1
\end{verbatim}

\section{Iloczyn tensorowy}

Dwa kubity opisujemy nie przez zwykłą parę wektorów, lecz przez iloczyn tensorowy przestrzeni stanów. Jeśli pierwszy kubit jest w stanie \(|a\rangle\), a drugi w stanie \(|b\rangle\), to stan złożony zapisujemy jako
\begin{equation}
    |a\rangle\otimes|b\rangle.
\end{equation}
Często skracamy zapis:
\begin{equation}
    |a\rangle|b\rangle=|ab\rangle.
\end{equation}
Baza dwóch kubitów ma cztery elementy:
\begin{equation}
    |00\rangle,
    |01\rangle,
    |10\rangle,
    |11\rangle.
\end{equation}

W reprezentacji macierzowej:
\begin{equation}
    |00\rangle=\begin{pmatrix}1\\0\\0\\0\end{pmatrix},
    \quad
    |01\rangle=\begin{pmatrix}0\\1\\0\\0\end{pmatrix},
    \quad
    |10\rangle=\begin{pmatrix}0\\0\\1\\0\end{pmatrix},
    \quad
    |11\rangle=\begin{pmatrix}0\\0\\0\\1\end{pmatrix}.
\end{equation}

\section{Stany dwukubitowe}

Ogólny stan dwóch kubitów ma postać
\begin{equation}
    |\psi\rangle
    =c_{00}|00\rangle+c_{01}|01\rangle+c_{10}|10\rangle+c_{11}|11\rangle.
\end{equation}
Warunek normalizacji to
\begin{equation}
    |c_{00}|^2+|c_{01}|^2+|c_{10}|^2+|c_{11}|^2=1.
\end{equation}
Jeżeli stan można zapisać jako iloczyn dwóch stanów pojedynczych kubitów, nazywamy go stanem separowalnym:
\begin{equation}
    |\psi\rangle=|a\rangle\otimes|b\rangle.
\end{equation}
Nie każdy stan dwóch kubitów ma taką postać. To prowadzi do pojęcia splątania.

\section{Splątanie}

Stan splątany to stan złożony, którego nie da się zapisać jako iloczynu stanów podukładów. Splątanie nie oznacza zwykłej korelacji wynikającej z niewiedzy. Oznacza, że opis całości jest bardziej podstawowy niż opis części osobno.

Najprostszy przykład to stan
\begin{equation}
    |\Phi^+\rangle=\frac{1}{\sqrt{2}}\left(|00\rangle+|11\rangle\right).
\end{equation}
Nie da się go zapisać jako
\begin{equation}
    (\alpha|0\rangle+\beta|1\rangle)\otimes(\gamma|0\rangle+\delta|1\rangle)
\end{equation}
przy odpowiednim doborze \(\alpha,\beta,\gamma,\delta\). Korelacje pomiarów dwóch kubitów są tutaj wbudowane w stan.

\section{Stan Bella}

Cztery standardowe stany Bella to
\begin{align}
    |\Phi^+\rangle &=\frac{1}{\sqrt{2}}(|00\rangle+|11\rangle),\\
    |\Phi^-\rangle &=\frac{1}{\sqrt{2}}(|00\rangle-|11\rangle),\\
    |\Psi^+\rangle &=\frac{1}{\sqrt{2}}(|01\rangle+|10\rangle),\\
    |\Psi^-\rangle &=\frac{1}{\sqrt{2}}(|01\rangle-|10\rangle).
\end{align}
Są to maksymalnie splątane stany dwóch kubitów. Tworzą ortonormalną bazę przestrzeni dwukubitowej.

W stanie \(|\Phi^+\rangle\), jeśli zmierzymy pierwszy kubit w bazie standardowej i otrzymamy \(0\), drugi też będzie \(0\). Jeśli otrzymamy \(1\), drugi też będzie \(1\). Przed pomiarem nie przypisujemy jednak każdemu kubitowi ukrytej klasycznej wartości. Stan opisuje korelację kwantową całości.

\section{Pomiar częściowy}

Pomiar jednego kubitu w stanie dwukubitowym zmienia opis całego układu. Dla stanu
\begin{equation}
    |\Phi^+\rangle=\frac{1}{\sqrt{2}}(|00\rangle+|11\rangle)
\end{equation}
pomiar pierwszego kubitu daje \(0\) z prawdopodobieństwem \(1/2\) i \(1\) z prawdopodobieństwem \(1/2\). Jeśli wynik to \(0\), stan po pomiarze jest zgodny z \(|00\rangle\). Jeśli wynik to \(1\), stan po pomiarze jest zgodny z \(|11\rangle\).

W formalnym opisie pomiar częściowy wymaga projektorów działających tylko na jeden czynnik iloczynu tensorowego, na przykład
\begin{equation}
    P_0\otimes I,
\end{equation}
gdzie
\begin{equation}
    P_0=|0\rangle\langle0|.
\end{equation}

\section{Interpretacja: czego splątanie nie oznacza?}

Splątanie nie oznacza wysyłania sygnału szybszego od światła. Nie oznacza też, że pomiar pierwszego kubitu klasycznie „popycha” drugi kubit przez przestrzeń. Oznacza, że stan złożony nie jest iloczynem stanów części, a wyniki pomiarów wykazują korelacje, których nie da się wyjaśnić prostym klasycznym modelem lokalnych ukrytych wartości.

Splątanie jest zasobem informatyki kwantowej, ale nie jest magiczną komunikacją. Żeby porównać wyniki pomiarów dwóch stron, nadal potrzebny jest klasyczny kanał komunikacji. Kwantowość leży w strukturze korelacji i w niemożliwości opisania całości jako zwykłej mieszaniny niezależnych stanów części.

\begin{summarybox}
Kubit pokazuje mechanikę kwantową w najczystszej postaci: stan jest wektorem w dwuwymiarowej przestrzeni Hilberta, operacje są macierzami unitarnymi, pomiar daje wyniki probabilistyczne, a układy złożone wymagają iloczynu tensorowego. Splątanie pojawia się wtedy, gdy stan całości nie rozkłada się na stany części.
\end{summarybox}
